\chapter{أنواع العصبونات}
\label{sec:neurontypes}

\section{العصبون صندوقًا أسود}

لنفترض أنك أُعطيت كيسًا فيه ستة وثمانون مليار عصبون~\cite{Azevedo2009}، وطُلب منك أن توزعها أكوامًا. كيف ستفعل ذلك؟

المهندس الذي يُعطى كيسًا من دوائر متكاملة لا يعرفها لن يفرزها بحسب شكل الغلاف. بل سيسأل: كم مدخلًا للقطعة، وكم مخرجًا، وماذا تُخرج. فلنرسم العصبون بالطريقة نفسها: كتلةً في مخطط (الشكل~\ref{fig:blackbox}).

\begin{figure}[t]
\centering
\begin{tikzpicture}[>=Stealth,x=1cm,y=1cm]
% the box
\node[draw,very thick,minimum width=2.6cm,minimum height=2.6cm,fill=blue!6] (N) at (0,0) {};
\node at (0,0.55) {\small عصبون};
\node at (0,-0.15) {$\displaystyle u=\sum_i w_i x_i$};
\node at (0,-0.8) {\small $y=\Theta(u-\theta)$};
% inputs
\foreach \k/\y/\lab in {1/1.05/{x_1},2/0.55/{x_2},3/0.05/{x_3}}{
  \draw[->,thick,blue!60!black] (-4.2,\y) -- (-1.3,\y);
  \node[left] at (-4.2,\y) {$\lab$};
  \node[above,blue!60!black] at (-2.1,\y-0.06) {\scriptsize $w_{\k}$};}
\node at (-4.45,-0.4) {$\vdots$};
\node at (-2.1,-0.4) {$\vdots$};
\draw[->,thick,blue!60!black] (-4.2,-1.05) -- (-1.3,-1.05);
\node[left] at (-4.2,-1.05) {$x_n$};
\node[above,blue!60!black] at (-2.1,-1.11) {\scriptsize $w_{n}$};
\draw[decorate,decoration={brace,amplitude=5pt},gray!60!black] (-4.9,-1.2) -- (-4.9,1.2);
\node[left,align=right,gray!40!black] at (-5.1,0) {\small $n\sim10^3\text{--}10^5$\\[-2pt]\small مدخل};
% single output wire, then fan-out
\draw[very thick,red!70!black] (1.3,0) -- (3.2,0);
\foreach \x in {1.6,2.2,2.34,2.9}{\draw[thick,red!70!black] (\x,0) -- (\x,0.4);}
\node[above right,font=\tiny\itshape,gray!30!black,inner sep=1pt] at (1.42,0.45) {نقرة\dots\ نقرة-نقرة\dots};
\node[below,red!70!black,align=center] at (2.25,-0.05) {\scriptsize مخرج واحد $y$};
\fill[red!70!black] (3.2,0) circle (0.06);
\foreach \y in {1.3,0.65,0,-0.65,-1.3}{
  \draw[->,thick,red!70!black] (3.2,0) -- (5.0,\y);}
\draw[decorate,decoration={brace,amplitude=5pt},gray!60!black] (5.3,1.45) -- (5.3,-1.45);
\node[right,align=left,gray!40!black] at (5.5,0) {\small نسخ من $y$\\[-2pt]\small إلى $10^3\text{--}10^4$\\[-2pt]\small عصبون آخر};
\end{tikzpicture}
\caption{العصبون بعين المهندس. مداخل كثيرة $x_i$، لكل منها وزنه $w_i$؛ وفي الداخل مجموع موزون $u$ ومقارنة بالعتبة $\theta$ ($\Theta$ دالة الدرجة: $1$ إذا كان المتغير موجبًا، و$0$ في غير ذلك)؛ ومخرج واحد $y$ هو سلسلة من النبضات تتكاثر بالتفرع لتصل إلى آلاف المستقبِلين.}
\label{fig:blackbox}
\end{figure}

ليس في مخرج الكتلة عدد، بل سلسلة من نبضات جهد قصيرة: ”نقرة“، ثم توقف، ثم ”نقرة-نقرة“، ثم توقف طويل. كل النبضات بالارتفاع نفسه، فالمعلومة كلها في \emph{توقيت} حدوثها. وفي داخل الكتلة، في تقريب أول خشن، جامع ومقارن: تُجمع المداخل بأوزانها، فإذا تجاوز المجموع العتبة أطلقت الكتلة نبضة. في الفصل التالي نستبدل بهذا النموذج نموذجًا يعمل بأمانة في الزمن المتصل.

المخرج واحد، لكن السلك يتفرع، وكل فرع يحمل \emph{النسخة نفسها} من الإشارة. يسمى هذا في الإلكترونيات التفرع عند المخرج (\foreignlanguage{english}{fan-out}). وهو عند عصبون القشرة من رتبة بضعة آلاف، بينما تحمّل البوابة المنطقية في الدائرة المتكاملة عادة ببضع بوابات أخرى فقط.

\emph{فما الذي} ينتقل عبر سلك المخرج إلى المستقبِل؟ تبلغ النبضة نهاية كل فرع، وهناك تتحول إلى إشارة كيميائية: يقذف الفرع في الشق الضيق بين الخليتين قليلًا من مادة معينة، هي \emph{الناقل العصبي}، فتغيّر تيار الدخل عند المستقبِل. وهذا هو معيار الفرز: \emph{أي ناقل عصبي يقذفه مخرج العصبون}.

\section{مخرج واحد، نوع واحد من الإشارة}

لا معنى للفرز إلا إذا كان لكل عصبون نوع واحد من المخرج. فهل الأمر كذلك؟ تقول التجربة إنه كذلك تقريبًا.

\begin{keyblock}
\begin{proposition}[عصبون واحد، نوع مخرج واحد]
\label{prop:dale}
لكل عصبون تقريبًا ناقل عصبي رئيسي واحد، ويقذفه في جميع فروع مخرجه. لذا يُحدَّد نوع العصبون بناقله الرئيسي~\cite{Strata1999}.
\end{proposition}
\end{keyblock}

لنتفق على أن نرمز إلى نوع العصبون \emph{بالأحرف الأربعة الأولى من الاسم الإنجليزي لناقله الرئيسي}. فالعصبون الذي ناقله الرئيسي الغلوتامات (\foreignlanguage{english}{glutamate}) هو عصبون \nt{GLUT}. الأنواع كلها مجموعة في الجدول~\ref{tab:types}.

\begin{table}[t]
\centering
\footnotesize
\setlength{\tabcolsep}{3pt}
\renewcommand{\arraystretch}{1.15}
\begin{tabular}{>{\raggedright}p{1.0cm}>{\raggedright}p{2.05cm}>{\raggedright}p{1.75cm}>{\raggedright}p{1.6cm}>{\centering}p{0.7cm}>{\raggedright}p{1.55cm}>{\raggedright}p{1.8cm}>{\raggedright\arraybackslash}p{3.2cm}}
\toprule
النوع & الناقل العصبي الرئيسي & الناقل (\foreignlanguage{english}{bus}) & السرعة & العلامة & العصبونات في الدماغ & مخارج العصبون & الدور في الحساب \\
\midrule
\nt{GLUT} & غلوتامات & عناوين & سريع وبطيء & $+$ & $\sim8\cdot10^{10}$ & $\sim10^4$ & الوزن الموجب الرئيسي \\
\nt{GABA} & غابا & عناوين & سريع وبطيء & $-$ & $\sim3\cdot10^{9}$ & $10^3\text{--}10^4$ & الوزن السالب الرئيسي \\
\nt{GLYC} & غليسين & عناوين & سريع & $-$ & $10^6\text{--}10^7$ & $10^2\text{--}10^3$ & وزن سالب في الحبل الشوكي وجذع الدماغ؛ ومفتاح ثانٍ لأحد مستقبلات الغلوتامات \\
\nt{ACET} & أسيتيل كولين & كلاهما & سريع وبطيء & $\pm$ & $\sim10^6$ & العضلة: $10\text{--}10^3$؛ الدماغ: $\sim10^5$ & مخرج إلى العضلة؛ وفي الدماغ سرعة التعلم (فرضية) \\
\nt{DOPA} & دوبامين & بثّ & بطيء & $\pm$ & $\sim5\cdot10^{5}$ & $10^5\text{--}10^6$ & خطأ التنبؤ بالمكافأة $\delta$ \\
\nt{SERO} & سيروتونين & بثّ & بطيء (نوع واحد من المستقبلات سريع) & $\pm$ & $\sim3\cdot10^{5}$ & $\sim10^5$ & أفق التخطيط (فرضية) \\
\nt{HIST} & هستامين & بثّ & بطيء & $+$ & $\sim6\cdot10^{4}$ & لا تقدير & إشارة عامة: ”ابقَ مستيقظًا“ \\
\nt{NORA} & نورأدرينالين & بثّ & بطيء & $\pm$ & $\sim5\cdot10^{4}$ & $\sim10^5$ & الكسب (فرضية) \\
\nt{ADRE} & أدرينالين & بثّ & بطيء & $\pm$ & $\sim10^4$ & لا تقدير & عصبونات قليلة؛ ودورها في الدماغ غير واضح \\
\nt{ASPA} & أسبارتات & عناوين & سريع & $+$ & لا تقدير & لا تقدير & وزن موجب ضعيف؛ ودوره موضع خلاف \\
\bottomrule
\end{tabular}
\caption{أنواع العصبونات، مرتبة تنازليًا بحسب عددها في الدماغ. \emph{الناقل}: ناقل العناوين، أي يُقذف الناقل العصبي في شق ضيق نحو مستقبِل محدد؛ وناقل البثّ، أي من مجموعة صغيرة من العصبونات المصدر إلى مناطق واسعة (القسم~\ref{sec:buses}). \emph{السرعة}: سريع، أي عبر المستقبلات مؤينة التوجه، في أجزاء من الألف من الثانية؛ وبطيء، أي عبر المستقبلات أيضية التوجه، في مئات من أجزاء الألف من الثانية وأكثر. \emph{العلامة} هي العلامة المعتادة؛ أما العلامة النهائية فيحددها المستقبِل (القسم~\ref{sec:receiver})، و$\pm$ تعني أن العلامتين كلتيهما موجودتان. \emph{العصبونات في الدماغ}: عدد عصبونات هذا النوع في دماغ الإنسان كله، من حيث رتبة المقدار؛ ومجموع العصبونات نحو $8.6\cdot10^{10}$~\cite{Azevedo2009}. معظم عصبونات \nt{GLUT}، أي نحو $7\cdot10^{10}$، هي عصبونات الدخل الصغيرة في المخيخ؛ وعددا \nt{GLUT} و\nt{GABA} مستخرجان من هذا العدّ ومن نسبة عصبونات \nt{GABA} في القشرة~\cite{Hendry1987}. ومن الأنواع القليلة العدد لم يُعدّ مباشرة إلا \nt{HIST}~\cite{Airaksinen1991} والمجموعة الرئيسية من مجموعتي عصبونات \nt{DOPA}~\cite{Pakkenberg1991}، لذا فالعدد في حالة \nt{DOPA} حدّ أدنى؛ وفي حالة \nt{SERO} هو مجموع عدّين جزئيين~\cite{Baker1990,Baker1991}. أما أعداد \nt{NORA} و\nt{GLYC} و\nt{ACET} و\nt{ADRE} فتقديرات خشنة لرتبة المقدار بلا عدّ مباشر. \emph{مخارج العصبون}: عدد النهايات المشبكية للعصبون الواحد، أي مواضع قذف الناقل العصبي على فروع سلك المخرج (الشكل~\ref{fig:blackbox})، من حيث رتبة المقدار. لم تُعدّ النهايات مباشرة في أنواع البثّ: فعددها عند \nt{DOPA} مستنتج من نسبة الهدف التي يغطيها تفرع سلك مخرج واحد~\cite{Matsuda2009}، وتقديرات البقية أخشن. ولـ\nt{ACET} حالتان: العصبون الذي يتحكم في عضلة، وعصبون البثّ في الدماغ.}
\label{tab:types}
\end{table}

% the pie is a table-type float with a figure caption, so it can never float ahead of Table~\ref{tab:types}
\begin{table}[tb]
\makeatletter\def\@captype{figure}\makeatother
\centering
\definecolor{vizglut}{HTML}{2A78D6}
\definecolor{vizgaba}{HTML}{E34948}
\definecolor{vizrest}{HTML}{8A8986}
\begin{tikzpicture}[>=Stealth]
% pie: GLUT 96.4 %, GABA 3.6 % (13.0 degrees), the rest 0.006 % (a hairline at 0 degrees)
\fill[vizglut] (0,0) -- (13:1.9) arc (13:360:1.9) -- cycle;
\fill[vizgaba] (0,0) -- (0:1.9) arc (0:13:1.9) -- cycle;
\draw[white,line width=1pt] (0,0) -- (13:1.9);
\draw[vizrest,line width=1.2pt] (0,0) -- (0:1.9);
\node[white,align=center,font=\small] at (190:0.95) {\nt{GLUT}\\$96.4\,\%$};
\draw[gray!60!black] (6.5:1.75) -- (2.05,2.1);
\node[anchor=west,font=\small] at (2.05,2.1) {\nt{GABA} $3.6\,\%$};
% zoom box for the remaining types
\draw[densely dashed,vizrest] (1.9,0) -- (3.1,1.65);
\draw[densely dashed,vizrest] (1.9,0) -- (3.1,-2.2);
\draw[vizrest,rounded corners=3pt] (3.1,-2.2) rectangle (10.1,1.65);
\node[anchor=west,font=\scriptsize] at (3.2,1.4) {كل الأنواع الأخرى معًا: $\approx0.006\,\%$؛ مقياس لوغاريتمي};
\foreach \code/\lg/\y/\val/\lx in {GLYC/6.477/1.0/{$10^6\text{--}10^7$}/4.05,ACET/6.0/0.6/{$\sim10^6$}/3.0,DOPA/5.699/0.2/{$\sim5\cdot10^5$}/2.699,SERO/5.477/-0.2/{$\sim3\cdot10^5$}/2.477,HIST/4.778/-0.6/{$\sim6\cdot10^4$}/1.778,NORA/4.699/-1.0/{$\sim5\cdot10^4$}/1.699,ADRE/4.0/-1.4/{$\sim10^4$}/1.0}{
  \node[anchor=east,font=\scriptsize] at (4.35,\y) {\nt{\code}};
  \fill[vizrest] (4.45,\y-0.13) rectangle ({4.45+\lg-3},\y+0.13);
  \node[anchor=west,font=\scriptsize] at ({4.5+\lx},\y) {\val};}
\draw[vizrest,thick] (7.45,1.0) -- (8.45,1.0);
\draw[vizrest,thick] (7.45,0.92) -- (7.45,1.08);
\draw[vizrest,thick] (8.45,0.92) -- (8.45,1.08);
\draw[gray!60!black] (4.45,-1.72) -- (8.45,-1.72);
\foreach \e in {3,4,5,6,7}{
  \draw[gray!60!black] ({4.45+\e-3},-1.72) -- ({4.45+\e-3},-1.66);
  \node[below,font=\scriptsize] at ({4.45+\e-3},-1.72) {$10^{\e}$};}
\end{tikzpicture}
\caption{حصص أنواع العصبونات في الدماغ بحسب تقديرات الجدول~\ref{tab:types}. تكاد عصبونات \nt{GLUT} تملأ الدائرة كلها؛ وعصبونات \nt{GABA} قطاع ضيق، أما كل الأنواع الأخرى معًا فشعرة على الحدّ. وعلى اليمين هذه الشعرة مكبّرة، بمقياس لوغاريتمي لعدد العصبونات؛ عمود \nt{GLYC} يبلغ منتصف المجال $10^6\text{--}10^7$، والقطعة تبيّن المجال كله. لا تقدير لـ\nt{ASPA}، فهو غائب عن الشكل.}
\label{fig:typepie}
\end{table}

يُظهر الشكل~\ref{fig:typepie} مدى التفاوت بين هذه الأكوام: من كل ألف عصبون $964$ من نوع \nt{GLUT}، و$36$ من نوع \nt{GABA}، أما الأنواع الأخرى كلها معًا فنصيبها نحو ستة عصبونات من كل مئة ألف.

اسم \foreignlanguage{english}{GABA} مؤلف أصلًا من أربعة أحرف. أما المواد التي لا تكاد تكون ناقلًا رئيسيًا، كالببتيدات العصبية والغازات والدهون والبيورينات، فلا تحصل على رمز خاص بها؛ والحديث عنها في الملحق~\ref{sec:othermediators}. كلمة ”تقريبًا“ في القضية~\ref{prop:dale} مهمة. فكثير من العصبونات يقذف مع الناقل الرئيسي موادَّ مساعدة~\cite{Yao2023}. وبعضها يقذف أيضًا ناقلًا سريعًا ثانيًا، قد يكون معاكسًا في العلامة: فجزء من عصبونات \nt{DOPA} يقذف الغابا أيضًا~\cite{Tritsch2012}. وفي بعض العصبونات تخرج نواقل مختلفة من فروع مختلفة للمخرج نفسه~\cite{Zhang2015}. كل هذا مُقاس، فقاعدة ”عصبون واحد، ناقل واحد“ قاعدة لها استثناءات وليست قانونًا. لكنها تصمد تقسيمًا أول: ففي الأطلس الخلوي لدماغ الفأر، حيث صُنّفت العصبونات بحسب الجينات العاملة فيها إلى أكثر من خمسة آلاف مجموعة، ما زالت العصبونات تنقسم إلى أصناف كبرى بحسب ناقلها الرئيسي أولًا~\cite{Yao2023}.

لهذه القاعدة نتيجة تميّز الدماغ فورًا عن الشبكة العصبية الاصطناعية. في الشبكة الاصطناعية يمكن للعصبون نفسه أن يكون له وزن موجب نحو مستقبِل ووزن سالب نحو مستقبِل آخر: فالأوزان $w_{ij}$ مجرد أعداد في مصفوفة. أما في الدماغ فعلامة التأثير يحددها في الغالب الناقل العصبي، وللعصبون ناقل واحد. إذن إذا أثار العصبون $j$ مستقبِلًا واحدًا فإنه يثير الباقين أيضًا. \emph{العمود} الخارج من العصبون $j$ في مصفوفة الأوزان كله بعلامة واحدة:
\begin{empheq}[box=\keybox]{equation}
\operatorname{sign} w_{ij} \;=\; s_j \quad\text{لكل المستقبِلين } i, \qquad s_j\in\{+1,-1\}.
\label{eq:dalesign}
\end{empheq}
تنقسم العصبونات إلى مثيرة ($s_j=+1$) ومثبطة ($s_j=-1$)، وهذه خاصية العصبون نفسه لا خاصية الوصلة. فإذا احتاجت الشبكة إلى وصلة سالبة من عصبون مثير، وجب أن تمر عبر وسيط: العصبون المثير يثير عصبونًا مثبطًا، وهذا يثبط الهدف، فينتج وزن سالب متأخر بحلقة واحدة.

النواقل العصبية المعروفة أكثر من مئة، لكن عمل الدماغ كله تقريبًا يقوم على نصف دزينة منها، وهي تنقسم إلى صنفين مختلفين تمامًا.

\section{ناقلان: ناقل العناوين وناقل البثّ}
\label{sec:buses}

في شبكات الحاسوب طريقتان لإيصال رسالة. الأولى \emph{عنوانية} (\foreignlanguage{english}{unicast}): تذهب الرزمة من مرسل واحد إلى مستقبِل محدد، بسرعة، ولا يراها أحد غيره. والثانية \emph{بثّ} (\foreignlanguage{english}{broadcast}): يرسل المرسل رسالة واحدة إلى كل عقد الشبكة دفعة واحدة. تستعمل العصبونات الطريقتين كلتيهما، وتنقسم النواقل العصبية بحسب هذا المعيار بالضبط (الشكل~\ref{fig:phone_radio}).

\begin{figure}[t]
\centering
\begin{tikzpicture}[>=Stealth,scale=0.95]
% ---- unicast: point-to-point wiring
\node[above] at (2.0,3.3) {\small \textbf{ناقل العناوين}: الغلوتامات والغابا};
\foreach \i in {0,...,4}{
  \fill[blue!60!black] (0.2+0.9*\i,2.5) circle (0.1);
  \fill[blue!60!black] (0.2+0.9*\i,0.6) circle (0.1);}
\foreach \a/\b in {0/0,0/1,1/1,1/3,2/2,3/2,3/4,4/4,2/0}{
  \draw[->,blue!60!black] (0.2+0.9*\a,2.38) -- (0.2+0.9*\b,0.72);}
\fill[red!70!black] (1.1,1.55) circle (0.09);
\pic[scale=1.2] at (1.75,1.9) {letter};
\draw[->,red!70!black,thick] (1.1,1.55) -- (0.28,0.7);
\draw[->,red!70!black,thick] (1.1,1.55) -- (1.95,0.72);
\node[align=center,below] at (2.0,0.2) {\small مليارات العصبونات،\\[-2pt]\small لكلٍّ منها مستقبِلوه،\\[-2pt]\small الزمن: أجزاء من الألف من الثانية};
% ---- broadcast: one small source, global targets
\begin{scope}[xshift=7.2cm]
\node[above] at (1.8,3.3) {\small \textbf{ناقل البثّ}: الدوبامين، \dots};
\foreach \i in {0,...,8}{\fill[blue!60!black] (-0.2+0.5*\i,2.75) circle (0.08);}
\fill[orange!85!black] (1.8,0.35) ellipse (0.35 and 0.18);
\pic[scale=1.5] at (2.7,0.75) {mast};
\foreach \x in {-0.2,0.3,0.8,1.3,1.8,2.3,2.8,3.3,3.8}{
  \draw[->,orange!85!black] (1.8,0.5) .. controls (1.8,1.3) and (\x,1.6) .. (\x,2.62);}
\node[align=center,below] at (1.8,0.1) {\small من آلاف إلى مئات الآلاف من العصبونات،\\[-2pt]\small رسالة واحدة للجميع،\\[-2pt]\small الزمن: مئات من أجزاء الألف من الثانية وأكثر};
\end{scope}
\end{tikzpicture}
\caption{طريقتان للإرسال. على اليسار ناقل العناوين، وهو يشبه البريد والعنوان مكتوب على الظرف؛ وبالأحمر عصبون واحد ومستقبِلاه. وعلى اليمين ناقل البثّ، كمحطة إذاعة: مجموعة صغيرة من العصبونات المصدر، والرسالة نفسها يتلقاها الجميع.}
\label{fig:phone_radio}
\end{figure}

\emph{ناقل العناوين} هو الغلوتامات والغابا. تقذفهما الغالبية الساحقة من العصبونات. كل مخرج يقود إلى خلايا محددة، ويعمل الناقل العصبي في الشق الضيق للتماس فقط، ويعمل بسرعة: في جزء من الألف من الثانية تنفتح القنوات الأيونية لدى المستقبِل، وفي بضعة أجزاء من الألف ينتهي كل شيء. هذا ناقل \emph{البيانات}: عبره ينتقل ما يحسبه الدماغ.

\emph{ناقل البثّ} هو الدوبامين والسيروتونين والنورأدرينالين، والأسيتيل كولين جزئيًا. تقذفها مجموعات ضئيلة من العصبونات، لكن مخرج كل منها يتفرع اتساعًا مذهلًا. وكثيرًا ما يُقذف الناقل العصبي لا في شق ضيق بل في الحيز بين الخلايا، فينتشر فيه ويؤثر في كل من حوله. ويعمل ببطء: لا يفتح القنوات مباشرة، بل يطلق داخل المستقبِل سلسلة من التفاعلات الكيميائية. لا تجري عبره البيانات، بل \emph{وسائط تحكم} مشتركة بين أجزاء كبيرة من الشبكة، تغيّر نمط عملها: الكسب، وسرعة التعلم، وإشارة ”كان هذا جيدًا“. لذلك تسمى هذه النواقل \emph{المعدِّلات}. ساد طويلًا الاعتقاد أن كل قناة من هذه القنوات عدد واحد للدماغ كله. وقد دققت القياسات الحديثة الصورة: القناة ليست سلكًا واحدًا، بل حزمة صغيرة من الخطوط المتوازية. فالعصبونات المصدر لناقل واحد تنقسم إلى أنظمة فرعية، لكل منها مستقبِلوه ورسالته، كما تضبط الإشارات المحلية كمية الناقل المقذوف في موضعه~\cite{Ren2018,Poe2020}. وعدد هذه الخطوط آحاد أو عشرات لا مليارات، فتبقى القناة قناة بثّ.

إن احتجت إلى صورة واحدة من هذا الفصل فهذه هي. الدماغ شبكة ضخمة تنقل البيانات بالعناوين، وفوقها بضع قنوات بثّ تحمل إشارات التحكم. سيعرف المبرمج هذا التصميم المألوف: حساب، ومعه مجموعة صغيرة من متغيرات التحكم، كل منها مشترك بين جزء كبير من الشبكة.

\section{\nt{GLUT}: الوزن الموجب}

الغلوتامات هو الناقل المثير الرئيسي في الدماغ. عندما يقذف مخرج العصبون الغلوتامات، تنفتح لدى المستقبِل قنوات يدخل عبرها الصوديوم، فيرتفع الجهد على غشاء المستقبِل، وتزداد فرصته لإطلاق نبضة خاصة به. وبلغة الشكل~\ref{fig:blackbox} هذا مدخل بوزن \emph{موجب}، $w_i>0$.

من يُخرج الغلوتامات؟ كل من ينقل البيانات إلى مسافات بعيدة تقريبًا: من ثلاثة أرباع إلى أربعة أخماس عصبونات القشرة، ومعظم العصبونات التي تربط مناطق الدماغ المختلفة. شبكة الدماغ هي أولًا شبكة من الوصلات الغلوتاماتية.

تفصيل هندسي واحد. بعد القذف يقفز تركيز الغلوتامات في الشق آلاف المرات، إلى نحو ميلي مول، ثم يهبط بثابت زمني نحو جزء من الألف من الثانية~\cite{Clements1992}: ينتشر الغلوتامات خارج الشق، وتضخه بروتينات خاصة عائدة به إلى الخلايا. لماذا؟ للسبب نفسه الذي يُعاد لأجله الإشارة إلى الصفر بعد كل رمز في خط الاتصال. إن لم يُزَل الناقل العصبي، فستصل النبضة التالية إلى خط ”مشغول“، وستندمج النبضات التي يفصل بينها بضعة أجزاء من الألف من الثانية.

\section{\nt{GABA}: الوزن السالب}

تخيّل شبكة كل أوزانها موجبة. إذا ولّدت كل نبضة في المتوسط أكثر من نبضة لاحقة واحدة، نمت الفعالية نموًا أسيًا واجتاحت الشبكة كلها في بضع خطوات. سيعرف مهندس الإلكترونيات هنا حلقة تغذية راجعة موجبة بكسب أكبر من الواحد: تذهب الدارة إلى الإشباع. ولتجنب ذلك نحتاج إلى أوزان سالبة. ومصدرها الرئيسي حمض غاما-أمينوبيوتيريك، واختصاره \foreignlanguage{english}{GABA} (غابا).

تفتح الغابا لدى المستقبِل قنوات تمرر أيونات الكلور. جهد اتزان الكلور قريب من جهد الراحة أو أدنى منه قليلًا، لذا تبقي قنوات الكلور المفتوحة الغشاءَ قرب الراحة ولا تدع الجهد يرتفع إلى العتبة. هذا مدخل بوزن \emph{سالب}، $w_i<0$. عصبونات \nt{GABA} تشكّل من خُمس إلى ربع عصبونات القشرة~\cite{Hendry1987}. مخارجها قصيرة، وهي تثبط جيرانها الأقربين.

التثبيط في الدماغ يفعل أكثر بكثير من مجرد الطرح. إنه يعمل تغذيةً راجعة سالبة تبقي الشبكة كلها عند نقطة التشغيل. ويُعتقد أن التيارين المثير والمثبط الواردين إلى العصبون في القشرة السليمة يكادان يتوازنان بدقة: هذا النمط تنبأت به النظرية~\cite{vanVreeswijk1996} وتُظهره قياسات التيارات في القشرة الحية~\cite{Haider2006}، فيبقى العصبون طوال الوقت قرب العتبة. والدارة المستقرة بتغذية راجعة سالبة قوية حول نقطة غير مستقرة تستجيب بسرعة وحساسية للإشارات الصغيرة؛ وهذه حيلة هندسية قديمة.

\section{\nt{DOPA}: إشارة الخطأ}

أول قناة بثّ هي الدوبامين. عدد العصبونات الدوبامينية عند الإنسان من رتبة نصف مليون، أي نحو عصبون واحد من كل مئة وسبعين ألفًا؛ هذا ما عُدّ في المجموعة الرئيسية من مجموعتيها~\cite{Pakkenberg1991}، فهو حدّ أدنى. تمتد مخارجها إلى المناطق التي تختار الأفعال.

ماذا تنقل هذه القناة؟ الجواب تعطيه تجارب تُسجَّل فيها نبضات العصبونات الدوبامينية لدى حيوان ينال مكافأة~\cite{Schultz1997}. يُعطى الحيوان من حين لآخر قطرة عصير على غير توقع، فتجيب العصبونات الدوبامينية برشقة قصيرة من النبضات. ثم يُضاء مصباح قبل العصير، دائمًا قبله بثانية. وعندما يتعلم الحيوان هذا الارتباط، تبدأ العصبونات بالاستجابة برشقة \emph{للمصباح}، ولا تستجيب إطلاقًا للعصير نفسه إذا جاء في موعده تمامًا. وإذا لم يُعطَ العصير بعد المصباح، \emph{تصمت} العصبونات في اللحظة التي كان ينبغي أن يأتي فيها، ويهبط ترددها دون المعتاد.

القاعدة التي تفسر الحالات الثلاث (الشكل~\ref{fig:rpe}): لا تخبر العصبونات بمدى جودة الأمر، بل بمدى كونه \emph{أفضل أو أسوأ مما كان متوقعًا}.

\begin{figure}[t]
\centering
\begin{tikzpicture}[>=Stealth,x=3.4cm,y=1cm]
\foreach \y/\lab in {2.8/{عصير غير متوقع},1.4/{عصير بعد المصباح},0/{مصباح بلا عصير}}{
  \draw[gray!70!black] (0,\y) -- (3,\y);
  \draw[densely dotted,gray!60!black] (0,\y+0.3) -- (3,\y+0.3);
  \node[left,align=right,font=\small] at (-0.03,\y+0.35) {\lab};}
% row 1: juice without a cue -> burst at the juice
\draw[very thick,orange!85!black] plot[domain=0:3,samples=120] (\x,{2.8+0.3+0.75*exp(-((\x-2.08)/0.07)^2)});
\pic[scale=0.8] at (2,2.58) {juice};
% row 2: cue then juice -> burst at the cue, nothing at the juice
\draw[very thick,orange!85!black] plot[domain=0:3,samples=120] (\x,{1.4+0.3+0.75*exp(-((\x-1.08)/0.07)^2)});
\pic[scale=0.85] at (1,1.15) {lamp}; \pic[scale=0.85] at (2,1.15) {juice};
% row 3: cue, juice omitted -> burst at the cue, dip when the juice was due
\draw[very thick,orange!85!black] plot[domain=0:3,samples=120] (\x,{0.3+0.75*exp(-((\x-1.08)/0.07)^2)-0.28*exp(-((\x-2.1)/0.12)^2)});
\pic[scale=0.85] at (1,-0.25) {lamp}; \pic[scale=0.85] at (2,-0.25) {nojuice};
\node[right,font=\scriptsize] at (2.36,0.1) {هبوط};
\node[right,font=\scriptsize] at (2.13,3.75) {مفاجأة!};
\node[left,font=\scriptsize] at (0.97,2.25) {دفقة عند المصباح};
\node[above,font=\scriptsize] at (2,1.75) {لا شيء};
\draw[->,gray!70!black] (0,-0.75) -- (3.05,-0.75) node[right,black,font=\small] {$t$};
\draw[gray!60!black] (1,-0.8) -- (1,-0.7) node[below=2pt,font=\scriptsize] {المصباح، $1\,\text{s}$};
\draw[gray!60!black] (2,-0.8) -- (2,-0.7) node[below=2pt,font=\scriptsize] {العصير، $2\,\text{s}$};
\end{tikzpicture}
\caption{تردد نبضات عصبونات \nt{DOPA} في ثلاث تجارب (تخطيطيًا، بحسب~\cite{Schultz1997}). الخط المنقط هو التردد الخلفي الثابت. المصباح هو الإشارة، والقطرة هي العصير، والقطرة المشطوبة عصير وُعد به ولم يُعطَ. الاستجابة هي خطأ التنبؤ~\eqref{eq:rpe}.}
\label{fig:rpe}
\end{figure}

المبرمج الذي يعرف التعلم المعزز~\cite{SuttonBarto2018} سيعرف هذه الكمية فورًا. الوكيل الذي يتعلم اختيار الأفعال يحتفظ بتقدير $V(s)$: مقدار المكافأة الذي يتوقعه وهو في الحالة $s$. وبعد كل خطوة يقارن التوقع بما ناله:
\begin{empheq}[box=\keybox]{equation}
\delta_t \;=\; r_t \;+\; \gamma\,V(s_{t+1}) \;-\; V(s_t) ,
\label{eq:rpe}
\end{empheq}
ثم يصحح التقدير: $V(s_t)\leftarrow V(s_t)+\alpha\,\delta_t$. هنا $r_t$ المكافأة المنالة، و$\gamma<1$ معامل الخصم (بكم تقل قيمة المكافأة المستقبلية عن الفورية)، و$\alpha$ سرعة التعلم. وتسمى الكمية $\delta_t$ \emph{خطأ الفروق الزمنية}.

وهي تسلك تمامًا سلوك العصبونات الدوبامينية. العصير غير المتوقع: $V$ ما زالت صفرًا، و$r>0$، و$\delta>0$. العصير المتعلَّم: تنبأ المصباح بالمكافأة، فـ$V(s_t)$ قبل العصير تساوي $r$ أصلًا، و$\delta=0$. العصير لم يأتِ: $V(s_t)=r$، لكن $r_t=0$، و$\delta<0$. أما الدفقة فتنتقل إلى المصباح لأن التوقع يقفز صعودًا في لحظة المصباح بالذات: $\gamma V(s_{t+1})-V(s_t)>0$. الخطوات $t,t+1,\dots$ هنا اصطلاح خاص بالخوارزمية: لا يوجد في الجسم مؤقت يعدّ الخطوات، وسنحصل على الكمية نفسها في الزمن المتصل في الفصل~\ref{sec:dofa}.

إذن الدوبامين بثٌّ لإشارة خطأ عددية $\delta$ إلى كل من يجب أن يتعلم بها. في تقريب أول: سلك واحد للدماغ كله وعدد واحد في كل لحظة؛ أما إلى أي حدّ هذا السلك في الواقع حزمة أسلاك فيناقَش في الفصل~\ref{sec:dofaalt}.

\section{قنوات البثّ الأخرى}

للمعدِّلات الأخرى فرضيات عمل فقط، تترجم أدوارها إلى اللغة نفسها لغة الوسائط العامة~\cite{Doya2002}. تعامل معها على أنها فرضيات تحديدًا.

\emph{\nt{NORA}، النورأدرينالين: الكسب.} العصبونات النورأدرينالينية أقل عددًا حتى من الدوبامينية، ببضع عشرات من الآلاف في تقدير خشن، ومخارجها تمتد إلى الدماغ كله تقريبًا. تنشط بحدة حين يحدث أمر غير متوقع أو مهم. يغيّر النورأدرينالين ميل منحنى النقل لدى العصبونات المستقبِلة: تضعف المداخل الضعيفة وتقوى القوية. ويُقاس ذلك هكذا: تُكبت النبضات الخلفية للمستقبِل أكثر مما تُكبت استجابته للمدخل، فتزداد نسبة الإشارة إلى الضجيج~\cite{Foote1975NE}. ويوصف هذا في نموذج بسيط بعدد واحد: إذا استجاب العصبون بتردد $f=F\bigl(g\cdot u\bigr)$ لتيار الدخل $u$، فإن النورأدرينالين يزيد المعامل $g$~\cite{AstonJones2005} (التفاصيل في الفصل~\ref{sec:nora}). الشبكة ذات $g$ الكبير تعمل بـ”تباين“ أعلى وتتخذ قراراتها أسرع؛ وذات $g$ الصغير تعمل ”بضبابية“ وتجرّب خيارات أكثر، كوكيل التعلم المعزز في نمط الاستكشاف.

\emph{\nt{ACET}، الأسيتيل كولين: سرعة التعلم.} يتحكم الأسيتيل كولين في مقدار ما تصغي فيه الشبكة إلى المدخل الحالي مقابل بياناتها المخزنة. عند ارتفاع مستواه تقوى المداخل الآتية من الحواس، وتضعف الوصلات الداخلية ”المستحضِرة“، ويسهل في الوقت نفسه تغيير الأوزان~\cite{Hasselmo2006}. بعبارة تقريبية: هو مفتاح ”كتابة/قراءة“، ومعه سرعة التعلم $\alpha$.

\emph{\nt{SERO}، السيروتونين: أفق التخطيط.} ما ينقله السيروتونين هو الأقل فهمًا. تربطه إحدى الفرضيات بمعامل الخصم $\gamma$ في~\eqref{eq:rpe}: المستوى العالي من السيروتونين يقابل الاستعداد لانتظار مكافأة كبيرة بدل انتزاع صغيرة الآن. تعمل العصبونات السيروتونينية بانتظام وببطء، ببضع نبضات في الثانية أثناء اليقظة، وتكاد تصمت في إحدى مراحل النوم~\cite{JacobsAzmitia1992}.

النواقل العصبية الستة كلها مجموعة في الجدول~\ref{tab:transmitters}.

\begin{table}[t]
\centering
\small
\begin{tabular}{>{\raggedright}p{2.4cm}>{\raggedright}p{3.0cm}>{\raggedright}p{2.0cm}>{\raggedright\arraybackslash}p{5.2cm}}
\toprule
نوع العصبون & حصة العصبونات & الناقل & الدور في الحساب \\
\midrule
\nt{GLUT} (غلوتامات) & الأغلبية & عناوين & وزن موجب، $w>0$ \\
\nt{GABA} (غابا) & $20\text{--}25\,\%$ في القشرة & عناوين & وزن سالب، $w<0$؛ تثبيت الاستقرار \\
\nt{DOPA} (دوبامين) & $\sim 6\cdot10^{-6}$ & بثّ & خطأ التنبؤ $\delta$ \\
\nt{NORA} (نورأدرينالين) & $\sim 6\cdot10^{-7}$ & بثّ & الكسب $g$ (فرضية) \\
\nt{ACET} (أسيتيل كولين) & صغيرة & كلاهما & سرعة التعلم $\alpha$؛ ”كتابة/قراءة“ (فرضية) \\
\nt{SERO} (سيروتونين) & $\sim 3\cdot10^{-6}$ & بثّ & الخصم $\gamma$ (فرضية) \\
\bottomrule
\end{tabular}
\caption{النواقل العصبية الرئيسية في الدماغ بلغة الحساب. العمود الأيمن تبسيط شديد: موثوق للغلوتامات والغابا والدوبامين، أما للبقية ففرضيات.}
\label{tab:transmitters}
\end{table}

\section{المستقبِل يحدد العلامة}
\label{sec:receiver}

لقد خدعتك قليلًا حين قلت: الغلوتامات وزن موجب، والغابا وزن سالب. لا تحمل أي جزيئة علامة في ذاتها. الجزيئة مجرد مفتاح. أما ما يحدث حين تُلتقط فيقرره \emph{القفل}، أي المستقبِل على الخلية المستقبِلة.

أفضل مثال هو الدوبامين. له مستقبلات من عائلتين~\cite{Beaulieu2011}. في إحداهما يسهّل إثارة الخلية، وفي الأخرى يصعّبها. وفي المنطقة التي تختار الأفعال تحمل بعض العصبونات مستقبلات من النوع الأول، وبعضها من النوع الثاني، فتقوّي إشارة الدوبامين نفسها مجموعةً وتضعف أخرى في آن واحد. الأمر أشبه برزمة بثّ واحدة تفسرها عقد الشبكة المختلفة تفسيرات مختلفة.

\begin{keyblock}
\begin{proposition}[المستقبِل يحدد معنى الرسالة]
\label{prop:receptor}
علامة تأثير الناقل العصبي وسرعته لا يحددهما الناقل نفسه، بل المستقبِل الذي يرتبط به. والمستقبلات صنفان. \emph{المستقبلات مؤينة التوجه} هي نفسها قنوات أيونية، وتنفتح في أجزاء من جزء الألف من الثانية: هذا تحكم مباشر في التيار، كبوابة الترانزستور. و\emph{المستقبلات أيضية التوجه} تطلق داخل الخلية سلسلة من التفاعلات الكيميائية وتعمل في مئات من أجزاء الألف من الثانية وأكثر: هذا أقرب إلى الكتابة في سجل إعدادات يغيّر عمل الخلية اللاحق. ناقل العناوين يتكلم في الغالب عبر الأولى، وناقل البثّ في الغالب عبر الثانية.
\end{proposition}
\end{keyblock}

فلماذا يصح أصلًا أن نقول ”الغلوتامات يثير“؟ لأن كل مستقبلات الغلوتامات التي نصادفها في الواقع تقريبًا مثيرة، وكل مستقبلات الغابا تقريبًا مثبطة. ليس هذا قانونًا من قوانين الطبيعة، بل اصطلاح موثوق جدًا، وعليه تقوم القاعدة~\eqref{eq:dalesign}.

\section{ما تحمله معك}

الغالبية الساحقة من العصبونات ناقلُ بياناتٍ بالعناوين، أوزان كل عصبون فيه بعلامة واحدة؛ وأقلية ضئيلة قنواتُ بثّ تنقل إلى أجزاء كبيرة من الدماغ بضع إشارات تحكم عامة. نحو عصبونين من كل مئة ألف، وعليهما يتوقف كيف تتعلم بقية الشبكة كلها وفي أي نمط تعمل. وهذا لا يدهش المهندس: ففي أي حاسوب سجلات التحكم أقل بما لا يقاس من خلايا الذاكرة، ومع ذلك لا تعمل الآلة من دونها.

في الفصل التالي سنختبر ما تستطيع أن تحسبه شبكة من عصبونات \nt{GLUT} وحدها، وهي أكثر الأنواع عددًا.

\section{تمارين}

\begin{exercise}[\lvlA]
\label{ex:01:1}
\emph{الأكوام والأسلاك.} بحسب الجدول~\ref{tab:types} (منتصف المجال بمقياس لوغاريتمي؛ \nt{ACET}: $10^5$ مخرج): (أ)~لو كان كل عصبون إنسانًا، فكم ضعفًا لسكان الأرض تبلغ عصبونات \nt{GLUT}، وبحجم أي مدينة عصبونات البثّ كلها؟ (ب)~كم ضعفًا تزيد حصة نهايات \nt{DOPA} و\nt{SERO} و\nt{NORA} و\nt{ACET} على حصة عصبوناتها؟
\end{exercise}

\begin{exercise}[\lvlA]
\label{ex:01:2}
\emph{العودة إلى الصفر.} يهبط الغلوتامات في الشق كـ$c_0e^{-t/\tau_c}$، حيث $c_0=1\mM$ و$\tau_c=1\ms$؛ ويلحظ المستقبِل التركيز إذا تجاوز $10\,\mu\text{M}$. (أ)~كم يبقى الخط ”مشغولًا“ بعد قذفة واحدة؟ (ب)~المضخات محجوبة جزئيًا، $\tau_c=10\ms$. حتى أي تردد للقذفات يلحق الخط أن يتحرر؟
\end{exercise}

\begin{exercise}[\lvlB]
\label{ex:01:3}
\emph{إلى أين تنتقل الدفقة.} في تجربة الشكل~\ref{fig:rpe} تتوالى بعد المصباح الحالات $s_0\to\dots\to s_{K-1}$ بخطوة $\Delta$، و$T=K\Delta$؛ تصل المكافأة $r$ في الانتقال الأخير، ولحظة المصباح غير متوقعة. أوجد قيم $V(s_k)$ المتعلَّمة واستجابة $\delta$ للمصباح. وبوضع $\gamma=e^{-\Delta/\tau_\gamma}$ أوجد $\tau_\gamma$ عند $T=1\,\mathrm{s}$ و$\Delta=0.1\,\mathrm{s}$ و$\gamma=0.95$.
\end{exercise}

\begin{exercise}[\lvlB]
\label{ex:01:4}
\emph{نمذجة: التعلم من الخطأ.} نمذج سلسلة التمرين~\ref{ex:01:3} مع $K=10$ و$\gamma=0.95$ و$r=1$ بالقاعدة $V(s_t)\leftarrow V(s_t)+\alpha\delta_t$. في أي محاولة $n^*$ تتجاوز الاستجابة للمصباح لأول مرة الاستجابة للعصير عند $\alpha=0.05$ و$0.1$ و$0.2$ و$0.5$؟ وكيف تتعلق $n^*$ بـ$\alpha$؟
\end{exercise}

\begin{exercise}[\lvlB]
\label{ex:01:5}
\emph{تصميم: بثّ عدد واحد.} يجب إيصال العدد $\delta$ إلى كل العصبونات الـ$10^{10}$؛ وكل مستقبِل، بما في ذلك المُرحِّلات، يجب أن يتلقى $10$ نهايات من الطبقة السابقة، وكل حلقة تضيف $1\ms$. كم عصبونًا وكم حلقة في أوفر شجرة مرحِّلات عند $10^4$ مخرج للعصبون (\nt{GLUT}) وعند $10^{5.5}$ (\nt{DOPA})؟
\end{exercise}

\begin{exercise}[\lvlB]
\label{ex:01:6}
\emph{لماذا التوازن حساس.} في الشبكة $80\,\%$ من \nt{GLUT} و$20\,\%$ من \nt{GABA}، وكل عصبون متصل بكل عصبون بأوزان $J_E/N$ و$-J_I/N$؛ $\tau\dot r=-r+Wr+h$. عند $J_E=10$ القيمة الذاتية الوحيدة غير الصفرية لـ$W$ تساوي $0.5$. أوجد نسبة التيار المثبط إلى المثير، وبكم في المئة يكفي إضعاف التثبيط لتنجرف الشبكة إلى انهيار متسلسل.
\end{exercise}

\begin{exercise}[\lvlC]
\label{ex:01:7}
\emph{بحث: هل \nt{SERO} هو $\gamma$؟} بحسب التمرين~\ref{ex:01:3} استجابة عصبونات \nt{DOPA} للمصباح تساوي $re^{-T/\tau_\gamma}$. التأخيرات $T=1,\dots,5\,\mathrm{s}$، ولكل منها $n$ عرضًا، ويُقاس $\ln\delta$ بتشتت $0.5$. كم $n$ يلزم للتمييز بين $\tau_\gamma=2\,\mathrm{s}$ و$2.4\,\mathrm{s}$ (المستوى $0.05$، القوة $0.8$)؟ وأي آلية غير $\gamma$ تعطي الهبوط نفسه مع $T$؟
\end{exercise}
